




%  A simple approach for this type of analysis is to aggregate the response variable by day and fit a time series model that incorporates daily covariates (for example, the proportions of respondents from specified geographic regions or the fraction of respondents that are internally displaced persons, or IDPs). However, there are many limitations to this aggregated analysis, not the least of which is the inability to accurately account for interactions among these factors at the household level. 

% Alternative methods for analyzing repeated surveys have long been available, with early developments explicitly modeling changes in parameters over time ~\citep{ScottSmith74,Scott77, Jones80, Tam} and further extensions developed by~\citet{Wong} and~\citet{Staudenmayer}. Multilevel models, first applied to repeated surveys by~\citet{DiPrete}, avoid the need for aggregation by fitting regression models with parameters that are permitted to vary over time. In other specifications, fixed effects may be held constant while a time series is introduced for daily random effects~\citep{Modugno}, and some more computationally intensive approaches account for autocorrelation in longer repeated surveys by combining multilevel modeling with an additional time series for small-scale temporal variation~\citep{Lebo}. When the primary objective is to estimate the statistical significance of model covariates throughout a time period of interest in a large and diverse geographic area, both spatial and temporal sources of variation must be considered. As such, we have developed a novel algorithm for the rigorous analysis of de-identified household level data that applies a hybrid modeling approach to account for regional and demographic variability in addition to temporal autocorrelation. Our methodology enables the fitting of easily interpretable regression models that fully leverage food security survey data without the loss of information inherent in aggregate analyses. 

% \section{How Additional White Noise Alters ARMA Processes}

To motivate challenging this approach, let's assume perfect estimation of fixed effect parameters ($\hat{\beta} = \beta$). In this unlikely event, we are able to take observed data and extract $\mathbf{w}:=X_t\uu + \bm{\varepsilon}$, which is a realization of an \textit{ARMA-Noise Model} \cite{WongMiller1990}. With no prior information about the latent time series, the best estimates for the elements of $\uu$ are the averages at each time step. These averages have additional noise that introduces bias in the estimation of the time series parameters. 

Being an AR(1) process, the covariance of $u_t$ and $u_{t\pm h}$ is given by Equation \ref{eq:acvf_ar}. The fact that $\frac{\gamma_{\uu + \mathbf{a}}(1)}{\gamma_{\uu + \mathbf{a}}(0)} \neq \frac{\gamma_{\uu + \mathbf{a}}(2)}{\gamma_{\uu + \mathbf{a}}(1)}$ precludes $\mathbf{\hat{\uu}} = \uu + \mathbf{a}$ from being an AR(1) process, but as we will now show, its ACVF is compatible with that of an ARMA(1,1). From Box (1976) \citep{box1976time}, we have the following theorem:

\textit{If $\{y_t\}$ follows an ARMA($p_1,q_1$) model, $\{a_t\}$ follows an ARMA($p_2,q_2$) and the two series are independent, then their sum $w_t = y_t + a_t$ follows an ARMA($p,q$) model in which $p \leq p_1 + p_2$ and $q \leq \text{max}(p_1 + q_2,p_2 + q_1)$. }

Because $\{a_t\}$ is white noise, $p_2 = q_2 = 0$. As $p_1=1$ and $q_1=0$, we can deduce that $\{w_t\}$ is an ARMA($p,q$) process where $p\leq 1$ and $q \leq 1$. Because $\gamma_{\hat{\uu}}(h+1)/\gamma_{\hat{\uu}}(h) = \phi $ for $h > 0$, we know that $p = 1$. The ACVF of $\hat{\uu}$ is therefore incompatible with that of an AR(1) process, so it must be ARMA(1,1). The fact that both ACVF's are altered by a factor of $\phi$ after lag 1 means that the AR component is unaffected by the addition of white noise. While noting the AR term $\phi$ goes unchanged by the added noise, the estimation thereof will be irredeemably biased if fit with an AR(1) model (See Figure \ref{fig:ar_bias}). $\xi_2 = (\phi,\sigma_\eta^2)$ must instead be estimated with a method that accounts for this additional noise


\begin{figure}
    \centering
    \includegraphics[width=.8\linewidth]{images/output/ar_bias.pdf}
    \caption{Average estimates for $\phi$ after fitting $\hat{\uu} = \uu + \mathbf{a}$ with an AR(1) model and an ARMA(1,1) model. The results were produced by 100 realizations with $T = 1000$, $\sigma_\eta^2 = \sigma_a^2$. The disparity between the accuracy of the two means grows with the relative magnitude of $\sigma_a^2$ with respect to $\sigma_\eta^2$.}
    \label{fig:ar_bias} 
\end{figure}
